\slajdid{10-s019}
\begin{frame}[label=10-s019]{Drugi način određivanja $V_{IH}$}
\gusttekst\setlength{\abovedisplayskip}{3pt}\setlength{\belowdisplayskip}{3pt}\setlength{\abovedisplayshortskip}{2pt}\setlength{\belowdisplayshortskip}{2pt}\[I_{B1}=\frac{V_I-V_{BE1}}{R_B}\qquad I_{C1}=\frac{V_{CC}-V_{CE1}}{R_C}\]
to je za tranzistor u zasićenju
\[I_{B1}=\frac{V_I-V_{BES}}{R_B}\qquad I_{C1}=\frac{V_{CC}-V_{CES}}{R_C}\qquad\beta_F I_{B1}>I_{C1}\]
Smanjenjem ulaznog napona smanjuje se struje $I_B$ i pri nekom ulaznom naponu $V_{I2}$ će se pojaviti
\[\beta_F I_{B1}=I_{C1}\]
\[\beta_F\frac{V_{I2}-V_{BES}}{R_B}=\frac{V_{CC}-V_{CES}}{R_C}\]
\[V_{I2}=\frac{R_B}{\beta_F R_C}(V_{CC}-V_{CES})+V_{BES}\]
i to je tačka $V_{IH}$. Ranije
\[V_{I1}=\frac{R_B}{\beta_F R_C}(V_{CC}-V_{CES})+V_{BE}\]
Uočite da se ova dva napona razlikuju
\[V_{I2}-V_{I1}=V_{BES}-V_{BE}\]
\end{frame}
